\chapter{Συμπεράσματα και μελλοντική εργασία}
\label{ch:conclusions}

\section{Συμπεράσματα}
\label{sec:concl-summary}

Η εργασία υλοποίησε δύο αλγορίθμους επιλογής \tl{top-k} σε τρεις
αρχιτεκτονικές, χρησιμοποιώντας κοινή υποδομή για την αξιολόγησή τους. Η
παρούσα ενότητα συνοψίζει τα συμπεράσματα του κεφαλαίου~\ref{ch:evaluation},
εστιάζοντας στα ευρήματα που διατηρούν την ισχύ τους ανεξαρτήτως του
υποκείμενου υλικού.

\paragraph{Κίνηση δεδομένων.} Για μικρό \tl{k} όλοι οι συνδυασμοί αλγορίθμων
και αρχιτεκτονικών περιορίζονται από το εύρος ζώνης και όχι από την
υπολογιστική ικανότητα των συσκευών (ενότητα~\ref{sec:eval-roofline}). Η
επίδοση εξαρτάται αποκλειστικά από τη μείωση του όγκου των δεδομένων που
μεταφέρονται. Αυτό επιβεβαιώνεται ξεκάθαρα από τη διτονική ταξινόμηση στη
μονάδα γραφικών. Αν και αξιοποιεί μεγαλύτερο μέρος του ταβανιού εύρους ζώνης
απ' ό,τι η βιβλιοθήκη αναφοράς, τελικά υστερεί επειδή αναγκάζεται να διακινήσει
πολλαπλάσια δεδομένα.

\paragraph{Ο λόγος $n/k$.} Η ταχύτερη υλοποίηση για \tl{k}=8 μετατρέπεται στη
βραδύτερη για \tl{k}=131072 ανεξαρτήτως πλατφόρμας
(ενότητα~\ref{sec:eval-time-k}). Η απόδοση του δικτύου ταξινόμησης παραμένει
πολύ πιο σταθερή. Ο αλγόριθμος του σωρού αρχίζει να φιλτράρει στοιχεία μόνο
αφού γεμίσει, περιορίζοντας τη διαθέσιμη παραλληλία βάσει του λόγου $n/k$.
Παρόμοια συμπεριφορά έχει παρατηρηθεί και στη
βιβλιογραφία~\cite{lin2026adaptivetopk}, αναδεικνύοντας την ανάγκη για
προσαρμοστική εναλλαγή αλγορίθμων ανάλογα με τα δεδομένα εισόδου.

\paragraph{Το όριο μέτρησης.} Ενώ η μονάδα γραφικών υπερτερεί στον καθαρό χρόνο
εκτέλεσης του πυρήνα για τα ακραία μεγέθη του πειράματος, η συνολική
\tl{end-to-end} μέτρηση αναδεικνύει τον επεξεργαστή ως ταχύτερο σε όλες τις
περιπτώσεις. Η μετατόπιση του ορίου μέτρησης αλλάζει πλήρως την τελική
κατάταξη. Η αξιολόγηση μιας υλοποίησης οφείλει να συνυπολογίζει το κόστος
μεταφοράς δεδομένων, διαφορετικά η σύγκριση ως προς το αν η εκφόρτωση στην
κάρτα συμφέρει παραμένει ελλιπής.

\paragraph{Ισχύς και ενέργεια.} Η μονάδα νευρωνικής επεξεργασίας έχει τη
χαμηλότερη ισχύ απ' όλες τις υλοποιήσεις και ταυτόχρονα την υψηλότερη
κατανάλωση, καθώς η διάρκεια κυριαρχεί του γινομένου
(ενότητα~\ref{sec:eval-energy}).

Για την επιλογή \tl{top-k} με μικρές τιμές του \tl{k}, η εκφόρτωση του
υπολογισμού σε επιταχυντή αποδεικνύεται ενεργειακά και χρονικά ασύμφορη
συγκριτικά με την εκτέλεση \tl{map-reduce} στον επεξεργαστή. Αρχιτεκτονικές
όπως οι \tl{NPU} (ενότητα~\ref{sec:bg-arch-npu}) προϋποθέτουν υψηλή τοπική
επαναχρησιμοποίηση δεδομένων για να αποσβέσουν το κόστος μεταφοράς τους. Η
επιλογή \tl{top-k} δεν προσφέρει αυτή τη δυνατότητα, με αποτέλεσμα το
μεγαλύτερο μέρος του χρόνου να αναλώνεται στην προετοιμασία των δεδομένων,
αφήνοντας τη συσκευή σε αδράνεια (ενότητα~\ref{sec:npu-pipeline}).

Μοναδική εξαίρεση αποτελεί η διτονική ταξινόμηση στη μονάδα νευρωνικής
επεξεργασίας για αριθμούς μισής ακρίβειας, όπου και ξεπερνά την αντίστοιχη
υλοποίηση αναφοράς. Το πλεονέκτημα αποδίδεται στην ελλιπή βελτιστοποίηση της
βιβλιοθήκης του επεξεργαστή για τον συγκεκριμένο τύπο δεδομένων
(ενότητα~\ref{sec:eval-dtypes}).

\section{Περιορισμοί}
\label{sec:concl-limitations}

Η σύγκριση με την \tl{std::partial\_sort} στον επεξεργαστή ενέχει μια εγγενή
ασυμμετρία. Η βιβλιοθήκη λειτουργεί αυστηρά μονονηματικά, σε αντίθεση με τις
υλοποιήσεις της εργασίας που εκμεταλλεύονται την παραλληλία
(ενότητα~\ref{sec:cpu-reference}). Η καταγεγραμμένη επιτάχυνση προκύπτει
συνδυαστικά τόσο από τη βελτίωση του αλγορίθμου όσο και από τη χρήση πολλαπλών
νημάτων. Η επιλογή αυτή πηγάζει από την πρόθεση να συγκριθεί η νέα υλοποίηση με
τα άμεσα διαθέσιμα εργαλεία ενός προγραμματιστή, ωστόσο η διαφορά αυτή πρέπει
να λαμβάνεται υπόψη στην ανάγνωση των αποτελεσμάτων.

Επιπλέον, ο όγκος δεδομένων που καταμετρήθηκε για τις βιβλιοθήκες αναφοράς
συνιστά απλώς το κάτω φράγμα. Επειδή τα εσωτερικά τους περάσματα στη μνήμη δεν
είναι άμεσα παρατηρήσιμα, τα αναφερόμενα ποσοστά αξιοποίησης της οροφής στην
ενότητα~\ref{sec:eval-bandwidth} αποτελούν υποεκτιμήσεις.

Για τη μονάδα νευρωνικής επεξεργασίας δεν διατίθεται εγγενής βιβλιοθήκη
\tl{top-k} για άμεση σύγκριση. Η αντιπαραβολή πραγματοποιήθηκε έναντι του
\tl{host}, απαντώντας στο πρακτικό ερώτημα του κατά πόσο αξίζει η εκφόρτωση,
χωρίς να αξιολογεί τη μέγιστη θεωρητική απόδοση της ίδιας της συσκευής.
Αντίστοιχα, η ανάλυση φάσεων της ενότητας~\ref{sec:npu-pipeline} προέρχεται από
ξεχωριστές εκτελέσεις με ενεργοποιημένο το \tl{debug} αντί από την τυπική
σάρωση, οπότε τα ποσοστά της στηρίζονται σε μικρότερο δείγμα επαναλήψεων απ'
ό,τι οι υπόλοιπες μετρήσεις του παραρτήματος~\ref{app:reproducibility}.

Τέλος, ο πειραματικός χώρος περιορίστηκε σε ένα συγκεκριμένο σύστημα και
συγκεκριμένα εύρη τιμών. Αν και οι ποιοτικές τάσεις γύρω από τον λόγο $n/k$ και
το ρόλο του ορίου μέτρησης αναμένεται να γενικεύονται σε παρόμοια συστήματα, οι
ακριβείς αριθμητικοί παράγοντες παραμένουν στενά συνδεδεμένοι με τη
συγκεκριμένη δοκιμαστική πλατφόρμα.

\clearpage

\section{Μελλοντική εργασία}
\label{sec:concl-future}

Σε επίπεδο συστηματικών βελτιστοποιήσεων, η εκτέλεση πολλαπλών ταυτόχρονων
ερωτημάτων \tl{top-k} είναι η κύρια μελλοντική επέκταση, ειδικά για τη μονάδα
νευρωνικής επεξεργασίας. Εφόσον το κύριο εμπόδιο εντοπίζεται στην προετοιμασία
και αποστολή των δεδομένων, η ομαδοποίηση των κλήσεων θα βοηθούσε καθοριστικά
στην απόσβεση αυτού του κόστους. Τέτοια επαναλαμβανόμενα ερωτήματα κυριαρχούν
στις εφαρμογές της ενότητας~\ref{sec:bg-problem}. Στην ίδια κατεύθυνση, η
αξιοποίηση μηχανισμών άμεσης πρόσβασης της συσκευής στη μνήμη του συστήματος,
αποφεύγοντας τη ρητή αντιγραφή δεδομένων από τον \tl{host}, θα μπορούσε να
βελτιώσει σημαντικά τη συνολική καθυστέρηση. Η υλοποίηση μιας τέτοιας
δυνατότητας εξαρτάται σε μεγάλο βαθμό από τις προδιαγραφές και τις διεπαφές της
εκάστοτε πλατφόρμας.

Ως προς τις αλγοριθμικές επεκτάσεις, ένας μηχανισμός δυναμικής εναλλαγής
αλγορίθμου θα αξιοποιούσε άμεσα το εύρημα του λόγου $n/k$. Γνωρίζοντας πλέον
πώς ο λόγος επηρεάζει την απόδοση, μια προσαρμοστική υλοποίηση θα μπορούσε να
επιλέγει αυτόματα μεταξύ σωρού και δικτύου ταξινόμησης για να διασφαλίζει ομαλή
απόδοση σε όλο το φάσμα μεγεθών. Επιπλέον, για μεγαλύτερες τιμές του $k$ όπου η
επιλογή προσεγγίζει την πλήρη ταξινόμηση, η διαχείριση εισόδων που δεν χωρούν
στη μνήμη της συσκευής αποτελεί φυσική συνέχεια, μετατρέποντας τη διαχείριση
της ροής από ευθύνη της υποδομής σε οργανικό τμήμα του ίδιου του αλγορίθμου.

Τέλος, η υποδομή αξιολόγησης που αναπτύχθηκε στην παρούσα εργασία μπορεί να
αξιοποιηθεί μελλοντικά για τη μελέτη και άλλων υπολογιστικών φορτίων. Ο
διαχωρισμός του αλγοριθμικού χρόνου από τον χρόνο μεταφοράς, η ενιαία καταγραφή
ενέργειας στις διαφορετικές αρχιτεκτονικές και η ανάλυση μέσω του μοντέλου
οροφής (\tl{roofline}) προσφέρουν ένα επαναχρησιμοποιήσιμο πλαίσιο, ικανό να
εφαρμοστεί σε οποιοδήποτε άλλο πρόβλημα εξετάζεται για εκφόρτωση σε ετερογενές
υλικό.
